\documentclass[a4paper,12pt]{article}
\usepackage[utf8]{inputenc}
\usepackage[english]{babel} % Adds processing for some simple special characters.

% Packages for better page formatting and nice footnotes. %
\usepackage{fancyhdr} % Headers.
\usepackage{multicol} % Columns.
\usepackage[margin=0.75in]{geometry} % Margins.
\pagestyle{fancy}
\fancyhf{} % What does this do?
\lhead{Alexander S. Wheaton}
\rhead{\today}
\cfoot{\thepage}

% Title and author for title page. %
\title{AGN Host Properties}
\author{
    Alexander S. Wheaton\\
    School of Physics and Astronomy\\
    The University of Edinburgh\\
    s1572046@ed.ac.uk\break
}

% Utilities for fine-grained control over image insertion. %
\usepackage{graphicx} % Insert images.
\usepackage{float} % Images float in document environment.
\usepackage{wrapfig} % Image/tables in multicols with \wrapfigure or \wraptable.
\usepackage{caption} % Captions for single images.
\usepackage{subcaption} % Captions for simultaneous images.
\graphicspath{ {./img/} }

% Not sure what this is for. %
\usepackage{fullpage,epsf}
\usepackage{lipsum}

% Some utilities for scientific and mathematical writing. %
\usepackage{amsmath} % Nice matrices.
\usepackage{xfrac}   % Slant fractions and other utilities.
\usepackage{isotope} % Nice markup syntax for chemical symbols.
\usepackage{siunitx} % Formatting for numbers with SI units.
\DeclareSIUnit\parsec{pc}
\DeclareSIUnit\year{yr}
\DeclareSIUnit\lightyear{ly}
\DeclareSIUnit\erg{erg}
\DeclareSIUnit\ago{a}

% Utilities for manual kerning adjustment. %
\newcommand\K{\kern.05em} % Small amount of kerning.
\newcommand\KK{\kern.1em} % Medium amount of kerning.
\newcommand\KKK{\kern.2em} % Large amount of kerning.
\newcommand\KKKK{\kern.3em} % Very large amount of kerning.

% Bibliography and referencing style. Use style=draft for troubleshooting.
\usepackage[backend=bibtex,style=phys,sorting=nyt]{biblatex}
\addbibresource{references.bib}

\begin{document}

\thispagestyle{empty}                   % No numbers of title page
\epsfxsize=40mm                         % Size of crest
\begin{minipage}[b]{110mm}
    {\Huge\bf School of Physics\\ and Astronomy
    \vspace*{17mm}}
\end{minipage}
\hfill
\begin{minipage}[t]{40mm}
    \makebox[40mm]{
    \includegraphics[width=40mm]{crest.eps}}
\end{minipage}
\par\noindent                                           % Centre Title, and name
% \vspace*{2cm}
\begin{center}
    \Large\bf \Large\bf MPhys Project\\
    \Large\bf Astrophysics\\[10pt]                     % Change to MP/CP/Astro
    \LARGE\bf Bayesian Inference of Star Formation History in the Host Galaxies of Tidal Disruption Events
\end{center}
\vspace*{0.5cm}
\begin{center}
    \bf Alexander S. Wheaton\\
    \today
\end{center}
% \vspace*{5mm}

\begin{abstract}
  \noindent In this paper I review some of the basic properties of active
  galactic nuclei (AGN) with emphasis on their tendency to be observed in host
  galaxies which are ``post-starburst.'' I discuss the physics of a second
  nuclear phenomenon, tidal disruption events (TDEs), in quiescent galaxies. I
  then present the spectra of eight such galaxies at $z < 0.1$ known to have
  recently hosted a tidal disruption event, observed with the XSHOOTER
  instrument on the Very Large Telescope (VLT). I use the \textsc{Bagpipes}
  Python module to explore the effects of different star formation history
  (SFH) parametrisations on observed host spectra. By ``blind-fitting''
  simulated spectra with a known SFH, I identify assumptions about the form of
  SFH which must be made in order to reliably detect past starbust activity in
  a now quiescent galaxy. I apply these assumptions to the spectra of the
  aforementioned TDEs, and find that their SFHs share the statistical behaviour
  of AGN---that they occur preferentially post-starburst galaxies.
\end{abstract}

\vspace*{1cm}

\subsubsection*{Declaration}
\begin{quotation}
  I declare that this project and report is my own work.
\end{quotation}

\hspace*{1cm}
\noindent Signature:\hspace*{1cm}\includegraphics[width=6cm]{signature_repaired.eps}
\hfill
Date: \today

\vfill
\noindent{\bf Supervisor:} Professor A. Lawrence, FRSE, FRaS
\hfill
22 Weeks

\newpage
\thispagestyle{empty}
\section*{Personal Statement}\label{sec:personal_statement}

The motivation for this project began when I was given access to the spectra and photometry for eight galaxies which recently hosted tidal disruption
events, gathered using the XSHOOTER instrument on the Very Large Telescope.
After early discussions with my supervisor, Andy Lawrence, we determined that
the aim of my project would be to investigate the statistical behaviour of the
star formation in the host galaxies of tidal disruption events. To do this, I
would utilise the \textsc{Bagpipes} Python module, authored in 2018 by one of
the Royal Observatory's very own, to infer the star formation histories of the
TDE hosts from their spectra.

I spent the first four weeks of semester one conducting a broad literature
review on the the statistical properties of AGN, their relationship to
starburst activity, and the proposed mechanisms for this relationship. I also
spent time during this period researching the various spectral signatures of
star formation. Research on these two topics has made up of the bulk of the
theory section of this work.

I spent the next four weeks learning to use the \textsc{Bagpipes} module to
simulate galactic spectra for blind fitting. This was slow work, as the module
is very much still in development, and its documentation remains sparse. I
spent of lot of time just reading the \textsc{Bagpipes} source code to
understand the fitting mechanisms it employs. During this period, I also spent
some time reviewing Bayesian inference, and stellar population fitting.

During the last three weeks of semester one, I simulated many model galaxies
with star formation at various redshifts to investigate whether I could reduce
the dimensionality of the fit based on some absorption/emission metric. This
was ultimately unfruitful, and so the results have not been included in this
report, but are still available at \url{github.com/aswheaton/agn-host-project}.
During this period, I also generated and exchanged model galaxy spectra with
Philip Short for blind fitting.

Over the next six weeks of semester two, I devised various routines and priors
for starburst detection, and blind tested them on the mock spectra generated by
Philip. For enhanced computing power, I had to learn to utilise MPI
parallelisation and the Slurm job scheduler on the Cuillin computing cluster at
the Royal Observatory Edinburgh, although this sometimes proved more trouble
than it was worth. During weekly meetings with Andy and occasionally Philip, we
discussed blind fitting results (and failures). Each week, I returned to the
fitting process with new ideas about how to constrain priors without
overfitting.

For two weeks, I broke from fitting briefly to return to the literature on
various galactic properties. With these in mind, I began to successfully detect
starbursts in the blind fitted spectra. By the end of this time Phil had
received and reduced the last of the XSHOOTER targets, and I turned my blind
fitting methods towards these data. The remaining weeks of semester two were
spent doing brief analysis of the TDE fitting outcomes, some comparison with
previous analyses of these targets, and the preparation of the presentation and
report.

\newpage
\thispagestyle{empty}
\tableofcontents

\newpage
\setcounter{page}{1} % Set page number to 1
\section{Introduction}\label{sec:introduction}

The apparent connection between active galactic nuclei (AGN) and star formation
has long been established, if not fully explained. Evidence continues to mount
for some kind of causal link between these phenomena, be it galactic wind,
heating of the interstellar medium, or some as of yet undiscovered mechanism.
Enquiry on this front has been limited by cosmic timescales. Although we can
observe whether galaxies are actively emitting from their nuclei, observations
of the ``duty cycle'' switching are rare.\cite{Lawrence_1987} So too for star
formation. Although we can observe trends in cosmic star formation rates (SFRs)
by peering at galaxies of higher and higher redshift, we still only see each
galaxy in snapshot, and the signatures of recent star formation fade in just a
few \SI{}{\mega\year}. A remedy for this limitation was developed by Carnall et
al. 2018: the \textsc{Bagpipes} Python module, which employs Bayesian inference
to parametrise the most probable star formation history of an
\textit{individual} galaxy from its spectrum.\cite{Carnall_2018}
\textsc{Bagpipes} gives us new tools to probe the relationships between nuclear
activity and star formation---an essential component in our understanding of
galaxy evolution. With this software, I turn to a related nuclear phenomena,
the tidal disruption event (TDE). Using the \textsc{Bagpipes} module and new
spectral data on eight TDEs from the Very Large Telescope (VLT) to investigate
what features of their star formation history can be reliably inferred and if
they, too, tend to occur in galaxies with recent star formation epochs. By
doing so, we create a more complete picture of nuclear processes, their
dynamics and impacts on their host galaxies, and take a small step towards a
more unified theory of AGN, TDEs, and galactic evolution.

\section{Theory \& Motivation}\label{sec:theory_and_motivation}
\subsection{Active Galactic Nuclei}\label{sec:active_galactic_nuclei}

The luminosity of most galaxies is dominated by stellar, gas, and dust
emission. The stellar component is the aggregate of many stars at different
temperatures emitting thermally as near-perfect
blackbodies.\cite{Peterson_1997} The medium of interstellar gas emits thermally
when heated by the host galaxy's constituent stars, and creates both absorption
and emission features at the particular ionisation energies of common
interstellar gases and metals. Larger grains of interstellar dust absorb and re-
emit stellar light in the infrared region of the spectrum.\cite{Dyson_1997}


Differences between the spectra of elliptical and disk galaxies hinge mostly
upon the age of their constituent stars. Ellipticals or ``early-types'' are
characterised by red photometric colours associated with an older stellar
population. They emit very little light at wavelengths less than
\SI{4000}{\angstrom}, and because their interstellar mediums are mostly
comprised of hot ionised gas, their spectra feature mostly absorption lines
typical of G- and K-type stars, with no emission lines.\cite{Mo_2010}

Disk or ``late-type'' galaxies are characterised by younger, bluer stars,
neutral \isotope{HI} and molecular \isotope{H_2} hydrogen regions.Their spiral
arms are defined by \isotope{HII} regions and dust features. The emit
considerably more light at wavelengths below the \SI{4000}{\angstrom} break, in
the blue and near-UV. Interstellar gas ionised by hot, young stars in these
galaxies produces strong emission lines, such as the Balmer-series
($\mathrm{H\alpha}$, $\mathrm{H\beta}$, $\mathrm{H\gamma}$, and
$\mathrm{H\delta}$), \isotope{OII} and \isotope{OIII} lines.\cite{Mo_2010}

A fraction ($\approx \sfrac{1}{3}$) of galaxies, however, feature tightly
confined (\SI{10}{\kilo\parsec}) ``weak'' nuclear luminosity which is
\textit{not} consistent with this model.\cite{McClure_2019} A smaller fraction
($\approx \sfrac{1}{100}$) exhibit nuclear luminosity which is a significant
fraction of the host's stellar luminosity ($\approx 10\%$).\cite{Sparke_2000}
These galaxies are said to have nuclei which are ``active'' and are referred to
as active galactic nuclei (AGN).\cite{Peterson_1997}

Lower luminosity AGN fall into the broad categories of ``Seyferts'' and radio
galaxies.\footnote{Additionally, a common and very low luminosity class of
galaxies featuring low-ionisation nuclear emission regions (LINERs), the highly
variable luminosity classes of BL Lac objects (blazars) and optically violent
variables (OVVs), and narrow-line X-ray galaxies, which are Seyferts heavily
reddened and extinguished by dust within the host.} Type 1 and 2 Seyferts
typically occur in spiral galaxies, featuring broad and narrow emission lines,
respectively. Broad-line radio galaxies (BLRGs) and narrow-line radio galaxies
(NLRGs) are the elliptical analogues to Seyferts, but are additionally radio-
loud.\cite{Peterson_1997}

Some AGN have power outputs many times that of a typical host galaxy. In these
objects, luminosity from the nucleus completely obscures the host. These are
believed to be exceptionally rare, and as such are only detected at very high
redshifts. At such distances, they appear point-like or  ``quasi-stellar,'' and
are so dubbed quasars.\cite{Peterson_1997} In general, AGN exhibit several key
characteristics:

\begin{itemize}
    \item highly luminous and compact appearance
    \item continuum luminosity which is not explicable by thermal emission
    \item variable luminosity over short timescales of days to months
\end{itemize}

Unlike the peaked blackbody emission from stars, AGN radiate across the entire
electromagnetic spectrum, with a peak in the UV and secondary peaks in the near
infrared and X-ray. The generally accepted explanation for continuum spectral
emission of this magnitude is a hot accretion disk surrounding a supermassive
black hole (SMBH).\cite{Sparke_2000} As matter falls into the black hole, its
potential energy is converted to kinetic, dissipated by viscosity in the
accretion disk to thermal energy, and radiated. In this model, material in the
disc at a range of different temperatures accounts for continuum emission,
while Doppler shifting of light emitted by disc material travelling at a range
of velocities creates broad emission lines.\cite{Peterson_1997}

Emission from the central black hole is said to have different ``modes''
depending on its rate of mass accretion. In a ``steady'' mode, the SMBH is
constantly fed by an infall of matter.\cite{Romero_Canizales_2016} This infall
is believed to be fueled by stellar wind driving matter from the halo of the
host galaxy into the nucleus, where it is accreted and consumed. When
accreting, these SMBHs can be as luminous as an entire galaxy
(\textgreater\SI{e45}{\erg\per\second}), and it is this type of accretion which
fuels Seyferts, radio galaxies, and their more luminous cousins,
quasars.\cite{Kormendy_2013, Peterson_1997}

\subsection{The Starburst-AGN Connection}\label{sec:starburst_agn_connection}

A significant fraction of AGN exhibit signatures of recent, large-scale star
formation. These signatures include exceptionally blue photometric colours and
strong \isotope{HII}-region-type emission lines, presumably gas ionised by hot,
young O- and B-type stars. They also feature strong radio emissions, consistent
with a high number of recent supernovae remnants.\cite{Peterson_1997}

Early investigations found that both Seyfert and radio galaxy luminosity was
commonly related to the luminosity of the host galaxy, or in other words, to
total stellar mass formed in that galaxy. In the late 20th century, a
relationship emerged between spheroid or bulge luminosity and the luminosity of
the constituent AGN in Seyferts and radio galaxies.\cite{Lawrence_1987}. Large
scale surveys subsequently find a strongly positive correlation between the
central black hole masses (formed by accretion) and the bulge luminosity and
the stellar velocity dispersion (driven by star formation) in neighbouring
galaxies.\cite{Veilleux_2008}

Other indicators of this relationship abound: Seyferts have intermediate age
stars and red giants in their nuclei contributing to continuum luminosity---a
strong signature of recent star formation. Additionally, more luminous Seyferts
are found to have much younger stellar populations than less luminous
Seyferts.\cite{Veilleux_2008}

The age of a stellar population can be characterised by several spectral
indices. The $D_n(4000)$ index is defined as the ratio of flux from
\SIrange{3850}{3950}{\angstrom} to that from \SIrange{4000}{4100}{\angstrom}. A
high density of metal lines in this region means that young, hot stellar
populations which ionise these metals will have a low $D_n$ with respect to
older populations.\cite{Kauffmann_2003a} Studies by Kauffmann et al. 2003c
found that the high luminosity AGN have substantially younger stellar
populations (as characterised by $D_n$) than lower luminosity AGN, for fixed
values of surface mass density. Quiescent galaxies do not exhibit this
relationship, which suggests the most powerful AGN exist in galaxies which have
recently undergone a short burst of star formation.\cite{Kauffmann_2003c}

Is this phenomenon really the result of a starburst, though, or simply ongoing
star formation in the host galaxy? Another spectral index, the
$\mathrm{H\delta}$ absorption line, or $\mathrm{H\delta_A}$, increases strongly
in galaxies where star formation ended \textless \SI{1}{\giga\year} ago.
Galaxies with continuous formation occupy a narrow locus within the $D_n(4000)$/
$\mathrm{H\delta_A}$ plane, so the distance from this locus can be used to
determine the ``burstiness'' of recent star formation history (SFH) in an AGN
host. Analyses from Kaufmann et al. 2013c find AGN to indeed have a much wider
spread about this locus---indicating greater variability in their SFHs than
normal galaxies.\cite{Kauffmann_2003c}

All of these point to some causal relationship between these phenomena.
Proposed mechanisms vary, but most involve mechanical and/or radiative feedback
from the starburst/AGN, which shuts off the fuel supply to the AGN or quenches
star formation in the host galaxy. Mechanical energy from starbursts is
provided by wind from O- and B- type stars for approximately the first
\SI{6}{\mega\year}, and is then dominated by supernovae for approximately
\SI{34}{\mega\year}. AGN may supply mechanical energy as high velocity gas
outflow, which may be ejecta from the accretion disk or collimated jets
accelerated by magnetic fields emanating from the accretion disk. Radiation
pressure from very luminous starbursts or AGN may create Compton-heated wind or
ionise hydrogen in the interstellar medium. All of above may inhibit the
collapse of molecular clouds to form stars, or may shut off the supply of fuel
to the central black hole.\cite{Veilleux_2008}

\subsection{Tidal Disruption Events}\label{sec:tidal_disruption_events}

A different type of nuclear phenomena are tidal disruption events (TDEs). In
otherwise quiescent galaxies with low-level to no nuclear emission, spectra are
dominated by stellar emission, as described in
\S\ref{sec:active_galactic_nuclei}. Most galaxies, including quiescent ones,
are still believed to host one or more SMBHs at their centres.\cite{Rees_1988}
In these quiescent galaxies, it is possible for the central SMBH to be in an
``intermittent'' mode, wherein it is mostly starved of fuel, but occasionally
consumes cold dust clouds or individual stars.\cite{Holoien_2016a, Rees_1988,
Evans_1989}

If a star approaches a black hole of mass $>10^6M_\odot$ with a pericentre
sufficiently close to be subject to tidal forces, the two can not simply be
treated as point masses. Inside the tidal disruption radius $r_T$, orbital
kinetic energy is greater than the gravitational binding energy of the star. A
tidal bulge forms on the surface of the star, which rotates to face the black
hole, but out of phase with its procession around the centre of gravity. A
tidal torque acts on the star, and the distribution of velocities applied by
this torque to different regions of the now elongated star tear it apart in a
``tidal disruption event.''\cite{Rees_1988, Evans_1989}

The radiative (and therefore observational) result of this event depends on the
fate of the resulting stellar debris. Approximately half of the stellar
material escapes unbound on various hyperbolic orbits, but more tightly bound
debris may later return within the tidal disruption radius of the black hole on
elliptical orbits.\cite{Rees_1988} Over successive periods on bound orbits,
this material gradually accretes on to the surface of the black hole, creating
a short-lived ($<\SI{1}{\year}$) increase in nuclear luminosity as disk
viscosity converts the material's potential energy to thermal
radiation.\cite{Evans_1989}

Observational data supports a slightly longer luminosity flare, with several
distinct phases. For a 100-200 day period after peak bolometric luminosity,
emission from the TDE flare is ``soft,'' or dominated by thermal emission in
the UV/optical region, which originates from the accretion disk and a short-
lived corona of stellar ejecta. Approximately 500 days after peak UV/optical
luminosity, this transitions to a phase of ``hard'' emission, which is
dominated by a power-law spectrum of higher energy photons, before finally
returning to a quiescent mode.\cite{Wevers_2021} Observational data suggests
that some accretion disks are stable for up to 5-10 years after peak
luminosity.\cite{van_Velzen_2019}

For reasons similar to those reviewed in \S\ref{sec:starburst_agn_connection},
the precise dynamics of these phenomena are of particular relevance to galactic
evolution models. Like steadily luminous Seyferts and radio galaxies, feedback
from the accretion process may have consequences for the galactic star
formation rate, or for AGN quenching. Analytical and numerical models predict
the velocity of stellar ejecta to be of sufficient order to shock the
interstellar medium (ISM) and inhibit steady accretion which would otherwise be
inevitable in a galaxy hosting a SMBH.\cite{Evans_1989} As the ejecta are
braked by the ISM, their kinetic energy is thermalised, heating the ISM,
possibly enough to stop the collapse of molecular \isotope{H_2} clouds into
star forming nurseries.\cite{Rees_1988}

Like AGN, observed TDEs often appear in galaxies which bear the signatures of
recent star formation, but our ability to investigate the particular dynamics
between the related processes of AGN, TDEs, and star formation is limited by
the timescales on which these processes occur. Although we can peer at high
redshift galaxies to determine \textit{trends} in galactic evolution, we can
not observe the evolution of \textit{individual} galaxies in real time. In this
project, I develop methods for quantifying the star formation histories of
individual galaxies, and apply them to the host galaxies of eight TDEs.

\subsection{The XSHOOTER Data}\label{sec:xshooter_data}

% Spare table with RAs and Decs included.
% \begin{table}[H]
%   \begin{tabular}{l c c}
%     TDE         & $\alpha$     & $\delta$     & Morphology \\
%     \hline \hline
%     ASASSN-14li & 12:48:15.23  & +17:46:26.44 & Compact    \\
%     ASASSN-15oi & 20:39:09.18  & -30:45:20.10 & Unresolved \\
%     AT2018fyk   & 22:50:16.090 & -44:51:53.50 & Unresolved \\
%     AT2019ahk   & 07:00:11.546 & -66:02:24.14 & E+A        \\
%     AT2019azh   & 08:13:16.945 & +22:38:54.03 & E+A        \\
%     AT2019dsg   & 20:57:02.974 & +14:12:15.86 & Unresolved \\
%     AT2019qiz   & 04:46:37.880 & -10:13:34.90 & Unresolved \\
%     iPTF16fnl   & 00:29:57.010 & +32:53:37.24 & E+A        \\
%     \hline
%   \end{tabular}
%   \caption{
%   XSHOOTER targets and their coordinates, in right ascension and declination.
%   }
%   \label{tab:xshooter_data}
% \end{table}

\begin{figure}[h]
  \centering
  \includegraphics[width=0.3\textwidth]{ASASSN-14li_host_red}
  \includegraphics[width=0.3\textwidth]{AT2019qiz_host_red}
  \includegraphics[width=0.3\textwidth]{iPTF16fnl_host_red}
  \caption{ASASSN-14li, AT2019qiz, and iPTF16fnl host galaxies, NED.\cite{NED_1990}}
  \label{fig:xshooter_hosts}
\end{figure}

To do this, I utilised spectroscopic and photometric data on eight galaxies
which hosted a tidal disruption event between 2014 and 2020. Table
\ref{tab:xshooter_data} lists the targets, which were observed between
September and December of 2020, with the XSHOOTER instrument on the Very Large
Telescope (VLT). They are all of similar magnitude and moderate redshift, all
being located at $z<0.1$. At such distances, they occupy little real estate on
a CCD (see Fig. \ref{fig:xshooter_hosts}), but the XSHOOTER instrument provides
detailed, ground-based spectroscopy of the TDE hosts, with flux bins of width
$\Delta\lambda=\SI{0.2}{\angstrom}$ below \SI{10200}{\angstrom} and
$\Delta\lambda=\SI{0.6}{\angstrom}$ above \SI{10200}{\angstrom}. Some of the
spectra featured strong Telluric lines from the Earth's atmosphere, but the
errors on these fluxes can be exaggerated to ensure \textsc{Bagpipes} does not
overfit them, as described in \S\ref{sec:bagpipes_module}.

\begin{table}[H]
  \centering
  \begin{tabular}{l l c r}
    TDE         & Host Galaxy             & $m$ & $z$      \\
    \hline \hline
    ASASSN-14li & PGC 043234              & 15  & 0.0206   \\
    ASASSN-15oi & 2MASX J20390918-3045201 & 16  & 0.0484   \\
    AT2018fyk   & LCRS B224721.6-450748   & 17  & 0.06     \\
    AT2019ahk   & 2MASX J07001137-6602251 & 17  & 0.026211 \\
    AT2019azh   & KUG 0810+227            & 15  & 0.022    \\
    AT2019dsg   & 2MASX J20570298+1412165 & 15  & 0.0512   \\
    AT2019qiz   & 2MASX J04463790-1013349 & 15  & 0.01513  \\
    iPTF16fnl   & Markarian 950           & 16  & 0.018    \\
    \hline
  \end{tabular}
  \caption{
  XSHOOTER target TDEs, their host galaxies, magnitudes, and redshifts, which
  were observed on the Very Large Telescope (VLT) between September and
  December 2020.\cite{Zwicky_1975, Jose_2014, Holoien_2016a, Arcavi_2015,
  Holoien_2016b, Wevers_2021, Cacella_2019, Holoien_2019, van_Velzen_2019,
  Perez_Torres_2019, Seibert_2019, Gezari_2016, Blagorodnova_2017}
  }
  \label{tab:xshooter_data}
\end{table}

\subsection{The Bagpipes Module}\label{sec:bagpipes_module}

\begin{figure}[H]
  \includegraphics[width=\textwidth]{alex_model_20_sfh}
  \includegraphics[width=\textwidth]{alex_model_20_spec}
  \caption{Parameterised star formation history and simulated spectrum with \textsc{Bagpipes}.}
  \label{fig:bagpipes_example_simulation}
\end{figure}

\textsc{Bagpipes} is a module for Python 3, implemented by Carnall et al.
2018.\cite{Carnall_2018} Its moniker is an acronym for Bayesian Analyses of
Galaxies for Physical Inference and Parameter EStimation. \textsc{Bagpipes}
constructs model galaxies from a variety of parameters, such as the form of the
SFH, observed redshift, dust attenuation, and nebular emission. From this
model, \textsc{Bagpipes} generates model spectral energy distributions, as they
would be observed if they were real galaxies. The luminosity per unit rest-
frame wavelength $L_\lambda(\lambda)$, from sum of four ``ingredients'':

\begin{itemize}
  \item The SFH, $\mathrm{SFR}(t)$, which is a linear combination of one or
  more of the functional forms described below.
  \item Simple stellar population models, $\mathrm{SSP}(a,\lambda,Z)$ which
  give the spectral energy distribution for a total stellar mass of age $a$
  given by the SFH, on the basis of an initial mass function (IMF) and the
  metallicity of the stellar population or populations.
  \item The transmission function of the ionised ISM, $T^+(a,\lambda)$ which
  determines how starlight is absorbed and re-emitted by hot gas regions.
  \item Finally, the transmission function of the neutral ISM,
  $T^0(a,\lambda)$, which determines how starlight is absorbed and remitted by
  cold gas and dust regions.\cite{Carnall_2018}
\end{itemize}

\noindent These are combined to give the galactic rest-frame spectrum:

\begin{equation}\label{eq:luminosity_per_rframe_wavelength}
  L_\lambda(\lambda) =
  \sum_{j=1}^{N_c} \sum_{i=1}^{N_a}
  \mathrm{SFR}_j(t_i)
  \mathrm{SSP}(a_i,\lambda,Z_j)
  T^+(a_i,\lambda)
  T^0(a_i,\lambda)
  \Delta a_i
\end{equation}

\noindent Where $i$ runs over the number of stellar age bins defined by
\textsc{Bagpipes} and $j$ runs over the components of the
SFH.\cite{Carnall_2018} An example spectrum from a parametrised SFH is given in
Fig. \ref{fig:bagpipes_example_simulation}. The parametrised SFH, is a linear
combination of one or more of the following functional forms. An exponentially
decaying burst with form:

\begin{equation}\label{eq:tau_model}
  \mathrm{SFR}(t) \propto
  \begin{cases}
    \exp\left({-\frac{t-T_0}{\tau}}\right) & t > T_0, \\
    0 & t < T_0
  \end{cases}
\end{equation}

\noindent To which I will refer as the tau model. A delayed exponentially
decaying component:

\begin{equation}\label{eq:delayed_model}
  \mathrm{SFR}(t) \propto
  \begin{cases}
    (t-T_0)\exp\left({-\frac{t-T_0}{\tau}}\right) & t > T_0, \\
    0 & t < T_0
  \end{cases}
\end{equation}

\noindent To which I will refer as the delayed model. A lognormal form:

\begin{equation}\label{eq:lognormal_model}
  \mathrm{SFR}(t) \propto
  \frac{1}{t\sqrt{2\pi\tau^2}}
  \exp\left({-\frac{\left(\ln{t-T_0}\right)^2}{2\tau^2}}\right)
\end{equation}

\noindent Defined by Gladders et al. 2013, where $t$ is the elapsed time since
the Big Bang, $T_0$ is the logarithmic decay time, and $\tau$ sets rise and
decay timescale.\cite{Gladders_2013} \textsc{Bagpipes}, however, parametrises
this as $t_{max}$, the time since the big bang at the peak rate of star
formation and the full width at half maximum (FWHM) of the star formation
rate.\cite{Carnall_2019b} Finally, \textsc{Bagpipes} allows a more versatile
double power law parametrisation:

\begin{equation}\label{eq:dblplaw_model}
  \mathrm{SFR}(t)\propto
  \left[
  \left(\frac{t}{\tau}\right)^{\alpha} +
  \left(\frac{t}{\tau}\right)^{-\beta}
  \right]
  ^{-1}
\end{equation}

\noindent Which usefully decouples the rising $\alpha$ and falling $\beta$
slopes of the star formation, at the expense of an additional parameter in a
fit which is already of high dimensionality and of increased computational
resources.\cite{Carnall_2018} Each of the above parametrisations include a
parameter each for the total stellar mass $M_\star$ formed by the component and
the metallicity $Z$ of the stars formed. It is also possible to include
``bursts'' of fixed mass $M_\star$ such that:

\begin{equation}\label{eq:burst_model}
  \mathrm{SFR}(t)\propto
  M_\star \delta(t-T_0)
\end{equation}

\noindent As a matter of semantics, such ``bursts'' are not necessarily
synonymous with the type of ``starburst'' activity for which we are searching.
The starburst component may take one of any of these forms.

Global parameters for the entire model galaxy include the observed redshift
$z$, the velocity dispersion $\sigma_v$ of the stellar population,the logarithm
of the nebular ionisation rate $\log{U}$, and the extinction $A_V$ in
magnitudes due to dust in the galaxy, which contribute to the summation in Eq.
\ref{eq:luminosity_per_rframe_wavelength}. \textsc{Bagpipes} allow several
different dust attenuation curves, but for all fits in this project I have
utilised the curve defined by Calzetti et al. 2000.\cite{Calzetti_2000}

\begin{figure}
  \includegraphics[width=\textwidth]{host_hyz_specwerr_fit}
  \includegraphics[width=\textwidth]{host_hyz_specwerr_sfh}
  \caption{Observed (blue) spectrum, fitted (orange) spectrum, and inferred star formation history with \textsc{Bagpipes}.}
  \label{fig:bagpipes_example_fit}
\end{figure}

The real power of \textsc{Bagpipes}, however, is that it can also do the
inverse: take in an observed spectrum and infer the most likely star formation
history for that galaxy. Fig. \ref{fig:bagpipes_example_fit} shows an example
of one such fit, with standard deviations from the median posterior SFH shown
in grey. Bayes' Theorem for some new data $\mathcal{D}$ and hypothesis
$\mathcal{H}$ gives that the posterior probability that the model parameters
$\Theta$ describe the data according to the hypothesis is:

\begin{equation}\label{eq:bayes_theorem}
  P(\Theta|\mathcal{D},\mathcal{H}) =
  \frac
  {P(\mathcal{D}|\Theta,\mathcal{H})P(\Theta|\mathcal{H})}
  {P(\mathcal{D}|\mathcal{H})}
\end{equation}

\noindent Where $P(\Theta|\mathcal{H})$ is the prior probability distribution
of parameters in the model. In combination with the \textsc{MultiNest} linear
algebra sampling method implemented by Buchner et al. 2014 in
\textsc{PyMultiNest}, \textsc{Bagpipes} uses this method of inference
to determine the probability that a model galaxy spectrum (the hypothesis)
described by Eq. \ref{eq:luminosity_per_rframe_wavelength} is a fit for the
observed spectrum (the data), and produce a posterior distribution of SFHs and
their associated parameters.\cite{Carnall_2018, Buchner_2014}

\subsection{Stellar Population Fitting}\label{sec:stellar_population_fitting}

\begin{table}
  \centering
  \begin{tabular}{c c r r r}
    Spectral Type & $M/M_\odot$ & $\lambda_{peak}$     & MS Lifetime   & No. Density \\
    \hline \hline
    O5            &	40          & \SI{725}{\angstrom}  & \SI{1}{Myr}   & 1           \\
    B0            &	15          & \SI{1035}{\angstrom} & \SI{11}{Myr}  & 5           \\
    A0            &	3.5         & \SI{2898}{\angstrom} & \SI{440}{Myr} & 25          \\
    F0            &	1.7         & \SI{3864}{\angstrom} & \SI{2.7}{Gyr} & 63          \\
    G0            &	1.1         & \SI{4830}{\angstrom} & \SI{8}{Gyr}   & 100         \\
    K0            &	0.8         & \SI{5796}{\angstrom} & \SI{17}{Gyr}  & 630         \\
    M0            &	0.5         & \SI{8280}{\angstrom} & \SI{56}{Gyr}  & 800         \\
    \hline
  \end{tabular}
  \caption{
  Typical masses, temperatures, and main-sequence lifetimes for seven Harvard
  spectral subtypes. Number densities per \SI{10000}{\cubic\parsec} for each
  subtype's respective super-type are adapted from Glenn et al.
  2001.\cite{Glenn_2001}
  }
  \label{tab:star_lifetimes_densities}
\end{table}

Many of the challenges in fitting spectra to SFHs are due to the limitations of
stellar population fitting, whereby the total luminosity and the ratio of blue
to red stars constituent in a host galaxy is used to infer the age of stellar
formation components. The main sequence lifetime of a star is given
by\cite{Prialnik_2010}:

\begin{equation}
  T_{MS} \propto \left(\frac{M}{M_\odot}\right)^{-2}
\end{equation}

\noindent From this we obtain the typical main-sequence lifetimes for stars of
different Harvard spectral classifications, given in Table
\ref{tab:star_lifetimes_densities}. From this data, the limitations of stellar
population fitting are immediately apparent. Due to the inverse power law
relationship between mass and main-sequence lifetime, small reductions in the
initial mass (and therefore, main-sequence temperature) of stars result in
dramatic increases in main-sequence lifetime. While on the main-sequence, the
luminosity and colour of these long-lived stars varies
little.\cite{Prialnik_2010} If, for example, the hottest stars in a stellar
population are G-type stars, then stellar population fitting may reasonably
find that population to be anywhere between \SI{3}{\giga\year} and
\SI{11}{\giga\year} old. This kind of temporal resolution is comparable to the
age of the universe, but star formation rates are believed to vary on much
smaller timescales.\cite{Kauffmann_2003c} If star forming processes create new
stars with initial masses according to some initial mass function (IMF), then
star formation of sufficient total mass formed will always create some stars of
each spectral type. But once the population of A-type stars, with lifetimes on
the order of \SI{100}{\mega\year}, die out completely, it becomes increasingly
difficult to pinpoint the precise age of that formation epoch.

Indeed, observations by Glenn et al. 2001 reproduced in Table
\ref{tab:star_lifetimes_densities} show that our own Milky Way, despite being a
relatively young, blue galaxy, is nevertheless mostly populated by cool red
dwarves, whose lifetimes are many times that of the current age of the universe.

\section{Methods: Blind Fitting with \textsc{Bagpipes}}\label{sec:blind_fitting}

Because stellar population fitting is so sensitive to specified priors, I first
sought to determine what SFHs could reliably be fit by \textsc{Bagpipes} before
turning it on the XSHOOTER data. To do this, I performed ``blind-fitting''
trials, whereby I simulated observed spectra from a known, priori SFH and then
fit these mock spectra with \textsc{Bagpipes} to see what priors are necessary
to reproduce the true SFH.\footnote{For clarity, I will hereafter use the
dative ``priori'' to describe properties of the simulated galaxy, the
nominative ``prior(s)'' to describe the parameters and parameter ranges used to
blind-fit the simulated galaxy, and ``posterior'' to describe the fitted
parameters.}

Because model galaxy spectra are simulated, their flux values are exact and

don't have associated uncertainties. So before they can be fit, they need to

have noise added to them. The particular choice of uncertainty to add is
somewhat arbitrary, but I tried to mimic the signal to noise ratio $S/N$ of the
XSHOOTER data when doing this. At wavelengths below approximately
\SI{3500}{\angstrom}, noise in the XSHOOTER data is on the same order of
magnitude as the signal. Between approximately \SIrange{3500}{9000}{\angstrom},
noise is one or two orders of magnitude less than the signal, and at longer
wavelengths, the $S/N$ decreases considerably.

\begin{figure}
  \includegraphics[width=\textwidth]{tde_sn_hist}
  \caption{Histograms of the $S/N$ ratio of flux values for the eight XSHOOTER
  targets. Note that some very low $S/N$ ratios are due to deliberately
  exaggerated errors on Telluric lines, as described in \S\ref{sec:xshooter_data}.}
  \label{fig:sig_to_noise_histograms}
\end{figure}

To produce similar uncertainties, I simulate an error for each flux value $f$
from a Gaussian distribution with $\mu=f$ and $\sigma=0.1f$. This doesn't
perfectly reproduce the wavelength-dependent uncertainty of the XSHOOTER data,
but errors sampled from the wings of the Gaussian reproduce some of the very
high and very low $S/N$ in the XSHOOTER data while maintaining the statistical
behaviour illustrated in Fig. \ref{fig:sig_to_noise_histograms}: that most
errors are one or two orders of magnitude less than their associated data
points.

The minimum parameter space covered by a \textsc{Bagpipes} fit includes a
single function SFH and its associated parameters, a redshift, a velocity
dispersion, and dust attenuation: a 7-dimensional parameter space. More nuanced
models which account for nebular emission, birth cloud lifetimes, and multiple
SFH components quickly grow to be very high-dimensional and therefore
computationally expensive. My two component fits were 13- or 14-dimensional,
depending on whether or not a double power-law function was included, and each
run took many hours on the Cuillin computing cluster at the Royal Observatory
Edinburgh. For these reasons, it's essential to fit in a way that is
computationally efficient.

As an aside, it's import to consider whether real star formation histories are
really like this. As discussed in \S\ref{sec:active_galactic_nuclei} and
\S\ref{sec:tidal_disruption_events}, star formation histories do not evolve on
sufficiently short timescales to observe them directly, but we can measure
trends in galactic star formation over varying redshifts. Stellar population
dating of high redshift galaxies suggests that cosmic star formation peaks in
most galaxies sometime around $1.5 < z < 3$, a period sometimes referred to as
``cosmic noon,'' and that star formation is ``quenched'' after this
time.\cite{Feldmann_2016} So there are compelling reasons to believe that most
of a galaxy's stellar mass should be formed in some bulk formation in the
distant past, with possibly variable, low-level star formation since that
epoch. The simulated spectra I produced for blind-fitting all conform to this
paradigm, so the blind fits are really a test of how reliably this kind of SFH
can be reconstructed from a spectrum. If, however, true SFHs do not follow this
trend, fits performed with priors that reflect this paradigm will still tend to
find star formation histories that affirm it---a form of model bias. In the
next four sections, I detail four ``runs'' or sets of priors and related
algorithms, hereafter referred to as R1, R2, R3 and R4. In each run, I test
different priors and fitting routines to reconstruct a known star formation
history from its simulated spectrum, while balancing the need to minimise model
bias with computational austerity.

\subsection{R1: Iterative Fitting}\label{sec:r1}

\begin{table}[h]
  \centering
  \begin{tabular}{l l r r l}
    Model       & Parameter            & Symbol/Unit                  & Range          & Prior   \\
    \hline \hline
    Exponential & Age                  & $t(z)-T_0/\mathrm{Gyr}$      & (3.5, 10)      & Uniform \\
    Delayed     & Decay time           & $\tau/\mathrm{Gyr}$          & (0.1, 2)       & Uniform \\
    \hline
    Lognormal   & Peak time            & $t_{max}/\mathrm{Gyr}$       & (3.5, 10.5)    & Uniform \\
                & FWHM                 & $\mathrm{FWHM/Gyr}$          & (0.1, 5)       & Uniform \\
    \hline
    Global      & Mass formed          & $M_\star/M_\odot$            & (0, 15)        & Uniform \\
                & Metallicity          & $Z/Z_\odot$                  & (0, 2.5)       & Uniform \\
                & Birth cloud lifetime & $T_{BC}/\mathrm{Gyr}$        & (0.005, 0.015) & Uniform \\
                & Velocity dispersion  & $\sigma_v/\mathrm{kms^{-1}}$ & (150, 200)     & Uniform \\
                & Extinction           & $A_V$                        & (0, 1)         & Uniform \\
                & Redshift             & $z$                          & (0, 0.1)       & Uniform \\
    \hline
  \end{tabular}
  \caption{Prior parameters and their distributions, as applied in the R1 round of blind fits.}
  \label{tab:r1_priors}
\end{table}

\begin{table}[h]
  \centering
  \begin{tabular}{l l r r l}
    Model & Parameter   & Symbol/Unit             & Range           & Prior   \\
    \hline \hline
    Burst & Age         & $t(z)-T_0/\mathrm{Gyr}$ & (0, $t_{low}$)  & Uniform \\
          & Mass formed & $M_\star/M_\odot$       & (0, 15)         & Uniform \\
    \hline
  \end{tabular}
  \caption{Prior parameters and their distributions, as applied in the R1 round of blind fits.}
  \label{tab:r1_burst_priors}
\end{table}

As discussed in \S\ref{sec:bagpipes_module}, computational resources are a
serious consideration when devising priors for a \textsc{Bagpipes} fit. So for
my first fitting run, I tried to first establish the functional form of the
old, bulk star formation, and then determine whether or not a more recent burst
improved the fit. This advantageously limits the first three fits to only ten
dimensions, and runs only one subsequent 13-dimensional fit.

I first fit each of the functional forms given in Tab. \ref{tab:r1_priors}. To
determine which of these best fits the spectral data, I calculate the chi-
squared, goodness of fit value:

\begin{equation}
  \chi^2 = \sum_{i}\frac{(O_i - C_i)^2}{\sigma_i^2}
\end{equation}

\noindent For each fit using the simulated data $O_i$, their errors $\sigma_i$,
and model spectrum $C_i$ fit by \textsc{Bagpipes}. I select the functional form
with the best fit (lowest chi-squared value) and run a second fit with that
functional form in linear combination with a delta function with priors given
in Tab. \ref{tab:r1_burst_priors}. For this second component, I constrain the
age of the burst to be less than $t_{low}$, which is taken from the median
posterior age of the bulk component, i.e. we are only interested in detecting
bursts which occur \textit{after} the bulk star formation. I perform a final
chi-squared calculation for this fit, and compare it to the single component
fit. If the old+burst form reduces the chi-squared statistic, then I consider
this a positive burst detection. If not, I consider the single component fit to
be the best, and take there to be no burst in the SFH.

Unfortunately, this seldom detected priori bursts in the posterior. Upon
inspection, I found that when fitting the bulk component without a burst,
\textsc{Bagpipes} favours posterior parameters with a much younger bulk
formation than the priori SFH in order to account for blue stars formed in the
burst. Constraining the age of the burst to be younger than this biases the
second fit towards lower burst masses and young bulk formation. Since the burst
can not be sufficiently old to account for blue stars in the spectrum, the bulk
formation can not be weighted back towards an older stellar population. It also
dubiously assumes that the functional form with the best single component fit
will also be the best fit in linear combination with a burst, which is not
necessarily the case. For these reasons, I chose to abandon this approach.

\subsection{R2 \& R3: Wide Parameter Space with Two Components}\label{sec:r2_and_r3}

\begin{table}[h]
  \centering
  \begin{tabular}{l l r r l}
    Model       & Parameter            & Symbol/Unit                  & Range          & Prior    \\
    \hline \hline
    Exponential & Age                  & $t(z)-T_0/\mathrm{Gyr}$      & (3.5, 10)      & Uniform \\
    Delayed     & Decay time           & $\tau/\mathrm{Gyr}$          & (0.1, 2)       & Uniform \\
    \hline
    Lognormal   & Peak time            & $t_{max}/\mathrm{Gyr}$       & (3.5, 10.5)    & Uniform \\
                & FWHM                 & $\mathrm{FWHM/Gyr}$          & (0.1, 5)       & Uniform \\
    \hline
    DPL         & Rising slope         & $\alpha$                     & (0, 10)        & Uniform \\
                & Falling slope        & $\beta$                      & (0, 10)        & Uniform \\
                & Turnover             & $\tau$                       & (0, 15)        & Uniform \\
    \hline
    Global      & Mass formed          & $M_\star/M_\odot$            & (0, 15)        & Uniform \\
                & Metallicity          & $Z/Z_\odot$                  & (0, 2.5)       & Uniform \\
                & Birth cloud lifetime & $T_{BC}/\mathrm{Gyr}$        & (0.005, 0.015) & Uniform \\
                & Velocity dispersion  & $\sigma_v/\mathrm{kms^{-1}}$ & (150, 200)     & Uniform \\
                & Extinction           & $A_V$                        & (0, 1)         & Uniform \\
                & Redshift             & $z$                          & (0, 0.1)       & Uniform \\
    \hline
  \end{tabular}
  \caption{Prior parameters and their distributions, as applied in the R2, and R3 round of blind fits.}
  \label{tab:r2_r3_priors}
\end{table}

I decided that in order to avoid suppressing the burst mass with young bulk
formation, both the age of the burst and the bulk formation should be allowed
to vary over a wider prior parameter space. Without the need to constrain the
burst by the median age of the bulk formation, I was able to fit the three
functional forms from R1 \textit{or} the more versatile double power law
function (which has no posterior age distribution). In these runs, I also
utilised a more physically realistic, exponentially decaying burst function
instead of the burst described by Eq. \ref{eq:burst_model}. This burst is
allowed to vary over a much larger parameter space:

\begin{table}[h!]
  \centering
  \begin{tabular}{l l r r l}
    Model       & Parameter   & Symbol/Unit              & Range    & Prior   \\
    \hline \hline
    Exponential & Age         & $t(z)-T_0/\mathrm{Gyr}$  & (0, 10)  & Uniform \\ % F-star lifetime
                & Decay time  & $\tau/\mathrm{Gyr}$      & (0.1, 2) & Uniform \\
                & Mass formed & $M_\star/M_\odot$        & (0, 15)  & Uniform \\
                & Metallicity & $Z/Z_\odot$              & (0, 2.5) & Uniform \\
  \end{tabular}
  \caption{Prior parameters and their distributions for the burst component, as applied in the R2 and R3 round of blind fits.}
  \label{tab:r2_r3_burst_priors}
\end{table}

The only difference between R2 and R3 is that in R3, the prior redshift $z$ of
the galaxy is constrained to be in the range of $z_{obs}\pm 0.001$, where
$z_{obs}$ is the observational redshift. (For blind-fitting, the
``observational'' redshift was simply taken to be the priori redshift). R2
simply let this prior vary between 0 and 0.1, a range which included all the
priori galaxies and the XSHOOTER targets. R3 did produce marginally more
accurate posterior redshifts, but both runs still failed to reliably detect the
priori bursts. Rather, they continued to favour younger bulk formation than in
the priori galaxy and suppressed the mass of bursts in the posterior
distribution.

\subsection{R4: Selecting Physically Reasonable Priors}\label{sec:r4}

\begin{table}[h]
  \centering
  \begin{tabular}{l l r r l}
    Model       & Parameter            & Symbol/Unit                  & Range             & Prior   \\
    \hline \hline
    Exponential & Age                  & $t(z)-T_0/\mathrm{Gyr}$      & (7.5, 12.5)       & Uniform \\
    Delayed     & Decay time           & $\tau/\mathrm{Gyr}$          & (0.5, 2)          & Uniform \\
    \hline
    Lognormal   & Peak time            & $t_{max}/\mathrm{Gyr}$       & (1, 6)            & Uniform \\
                & FWHM                 & $2\sigma/\mathrm{Gyr}$       & (0.1, 3)          & Uniform \\
    \hline
    DPL         & Rising slope         & $\alpha$                     & (5, 10)           & Uniform \\
                & Falling slope        & $\beta$                      & (5, 10)           & Uniform \\
                & Turnover             & $\tau$                       & (0, 4.5)          & Uniform \\
    \hline
    Global      & Mass formed          & $M_\star/M_\odot$            & (5, 12.5)         & Uniform \\
                & Metallicity          & $Z/Z_\odot$                  & (0, 2.5)          & Uniform \\
                & Birth cloud lifetime & $T_{BC}/\mathrm{Gyr}$        & (0.013, 0.021)    & Uniform \\
                & Nebular ionisation   & $\log{U}$                    & (-4, -2)          & Uniform \\
                & Velocity dispersion  & $\sigma_v/\mathrm{kms^{-1}}$ & (50, 450)         & Uniform \\
                & Extinction           & $A_V$                        & (0, 2)            & Uniform \\
                & Redshift             & $z$                          & $z_{obs}\pm0.001$ & Uniform \\
    \hline
  \end{tabular}
  \caption{Prior parameters and their distributions, as applied in the R4 round of blind fits.}
  \label{tab:r4_priors}
\end{table}

\begin{table}[h]
  \centering
  \begin{tabular}{l l r r l}
    Model       & Parameter   & Symbol/Unit             & Range      & Prior   \\
    \hline \hline
    Exponential & Age         & $t(z)-T_0/\mathrm{Gyr}$ & (0.0, 3.5) & Uniform \\ % F-star lifetime
                & Decay time  & $\tau/\mathrm{Gyr}$     & (0.1, 2)   & Uniform \\
                & Mass formed & $M_\star/M_\odot$       & (0, 12.5)  & Uniform \\
                & Metallicity & $Z/Z_\odot$             & (0, 2.5)   & Uniform \\
    \hline
  \end{tabular}
  \caption{Prior parameters and their distributions on the burst component, as applied in the R4 round of blind fits.}
  \label{tab:r4_burst_priors}
\end{table}

I concluded from R2 and R3 that, actually, the age of the bulk formation had to
be constrained to some range more consistent with observation in order to
detect recent bursts, otherwise the burst will be ``smeared out'' and lost in
some middle aged star formation component. Most galaxies undergo quenching
sometime after cosmic noon.\cite{Tomczak_2014, Feldmann_2016} To reflect this
in my priors, I constrained the peak rate of bulk star formation to occur no
later than \SI{6.5}{\giga\year} after the Big Bang, and put a lower bound on
the stellar mass formed by this component. The reasoning for this boils down to
a few basic ideas:

\begin{itemize}
  \item The precise functional form of the old star formation is hard to pin
  down, due to the limits of stellar population fitting. As discussed in
  \S\ref{sec:stellar_population_fitting}, temporal resolution of a population
  decreases dramatically after its O- and B-type star die out.
  \item For the same reason, though, the particular form of the old component
  is perhaps unimportant for the purpose of starburst detection, for which one
  simply needs to establish whether or not there are two different stellar
  populations. For this purpose, the various functional forms of star formation
  will produce very similar stellar populations in the present if they are of
  similar mass and sufficient age, and whichever one gives the better spectrum
  fit is fine.
  \item This is physically reasonable and consistent with observation---high
  redshift observations find peak star formation around cosmic noon, and we do
  not observe newly formed galaxies at low redshifts, so most of a galaxy's
  mass should be formed in some older event, not in subsequent starbursts.
\end{itemize}

\noindent For the lower and upper mass bounds on this formation, I used the
standard deviations of the mean galactic mass for galaxies with $z \leq 0.1$,
which are $10^{5}M_\odot$ and $10^{12.5}M_\odot$.\cite{Baldry_2012,
Wright_2017} The burst component is constrained to occur no earlier than
\SI{3.5}{\giga\year} ago---slightly more than the typical main-sequence
lifetime for an F-type star. The mass formed in the burst is allowed to be as
large or larger than the old component, but \textit{there is no lower mass
limit for the burst.} The reasoning for this prior constraint is that, bursts
older than lifetime of these stars will probably be difficult to distinguish
from the low-level of star formation in quiescent galaxies, which for the Milky
Way is approximately $0.68-1.45M_\odot\mathrm{yr}^{-1}$, so there may be little
use in trying to fit them.\cite{Robitaille_2010} If the priori SFH does not
contain a burst, then the posterior mass of the burst should favour zero mass
formed. This gives a mechanism for a null result, while still exploring a wide
parameter space in a single run.

If the median posterior burst mass and decay time create a burst which exceeds
the quiescent star formation rate (or decaying rate from the old formation
component), I consider this a positive burst detection. If the median posterior
burst mass approaches or exceeds the total mass formed by the old component, it
is worth considering whether over fitting is at play, and the fit is trying to
compensate for an epoch of star formation between
\SIrange{6.5}{10.25}{\giga\year} after the Big Bang. This should be accountable
by the old component; although the \textit{peak} SFR of the old component is
constrained to be no later than \SI{6.5}{\giga\year}, the prior decay times/
falling slopes of these components allow star formation to continue well up
until the burst prior, even though they are quenching during this period.

During this time, I also went to the literature to constrain other parameters
to be consistent with observational values. The redshift $z$ varies tightly
around the observational redshift $z_{obs}$ within the range of $z_{obs} \pm
0.001$. Velocity dispersion $\sigma_{v}$ of constituent stars is limited by the
Faber-Jackson relation, and is therefore constrained to be
\SIrange{50}{450}{\kilo\meter\per\second}.\cite{Minkowski_1962} The prior on
the lifetime of the stellar birth clouds was chosen to be a normal distribution
about ${t_{BC}=17Myr}$ with ${\sigma=4Myr}$, consistent with typical cloud
lifetimes in the Milky Way, as determined by Murray 2011.\cite{Murray_2011} The
ionization parameter for nebular emission is allowed to vary between $\log{U}=-
4$ and $\log{U}=-2$, a very low and very high rate of gas ionization. The prior
dust component uses a Calzetti extinction curve, with $A_V$ allowed to vary
between one and two magnitudes (which may be on the optimistic side).

% \footnote{Note that the delayed form was actually run with $0.1 < \tau <
% 2.0$. This was in error, but doesn't seem to have seriously affected the
% posterior results.}

\section{Blind Fitting Results}\label{sec:blind_fitting_results}

With these priors, I ran a final round of blind fits, and finally began to
detect starbursts in the posterior. Each run has two components: one of the
four functional forms from Tab. \ref{tab:r4_priors} and the burst component
from Tab. \ref{tab:r4_burst_priors}. Of the four runs, I select the one with
the best chi-squared statistic, and investigate its posterior values to
determine whether a burst has been positively detected. Several blind fit
results are given below, each of which illustrates a different priori case or
posterior outcome, and the usefulness/limitations of this set of
priors.\footnote{A complete inventory of blind fitting results including model
spectra, fit priors, and posterior samples
in \texttt{.h5} format for all runs R1-R4 is available at \url{github.com/
aswheaton/agn-host-project/pipes}.}

\begin{figure}[h]
  \includegraphics[width=\textwidth]{../pipes/plots/r0_priors/phil_model_4_sfh}
  \includegraphics[width=\textwidth]{../pipes/plots/r4_lognormal_burst/phil_model_4_sfh}
  \caption{
  Blind fit model No. 4, priori SFH (top) and median posterior SFH (bottom)
  with $\chi^2=28.517106832593093$. Standard deviations of the posterior SFH
  are given in grey.
  }
  \label{fig:phil_model_4_sfhs}
\end{figure}

\begin{figure}[h]
  \includegraphics[width=\textwidth]{phil_model_4_posterior}
  \caption{
  Posterior distributions for fitted parameters on model No. 4. From left to
  right, the age of the burst, the mass formed by the burst, the peak of the
  old component, and the mass formed by the old component. Median and standard
  deviations of the distributions are shown by the dashed lines.
  }
  \label{fig:phil_model_4_posterior}
\end{figure}

Model No. 4 is an example case of a positive burst detection. Fig.
\ref{fig:phil_model_4_sfhs} shows the priori and posterior SFH for this model
galaxy. The posterior SFH contains a lognormal old component and a clear burst
component with median posterior age \SI{0.78}{\giga\year}. \textsc{Bagpipes}
assumes a flat $\lambda$-CDM cosmology such that the universe's present age is
\SI{13.467}{\giga\year}. At the median posterior redshift $z=0.069$ then, the
model galaxy's age is \SI{12.54}{\giga\year} and from Eq. \ref{eq:tau_model},
we have that $T_0 = \SI{11.78}{\giga\year}$. This is very close to the priori
$T_0=\SI{11.54}{\giga\year}$.

The stellar mass formed in the posterior burst is $M=10^{7.99}M_\odot$,
slightly less than the priori mass of $10^{8.3}M_\odot$. To understand why this
is the case, we consider also the posterior mass of the old component, which is
$10^{9.49}M_\odot$, compared to the priori mass (a double power law model) of
$10^{9.5}M_\odot$. These are very similar, so it is unlikely that mass has been
shifted from the burst to the old component. The time distribution of mass, on
the other hand, has shifted significantly, with the posterior old component peak
$t_{max} = \SI{5.55}{\giga\year}$ occurring much later than the priori turnover
time $\tau = \SI{4}{\giga\year}$. Additionally, the FWHM of the posterior old
component is much wider (\SI{2.81}{\giga\year}) than the width of the priori
double power law component, allowing star formation to continue long after its
peak. So stellar mass from old components can be smeared towards younger ages,
affecting the shape of the old components---the posterior SFR in Fig.
\ref{fig:phil_model_4_sfhs} peaks much lower than its priori counterpart. A
secondary effect of this is that some fraction of blue stars created in the
priori burst are still being accounted for by the decaying old component.
Clearly, even with this set of constraints on the priors, some mass of the
burst can be lost in the fit. Fig. \ref{fig:phil_model_4_posterior} shows the
1-dimensional posterior distributions of these key parameters. The distributions
are sharply peaked with small standard deviations, so the fit gives a high
level of confidence in these values, even though they are not consistent with
the priori model. In the future, this might be solved by further constraining
the decay times of the old component to quench more rapidly.

% parameter     median     16th percentile     84th percentile
% dust:Av:  0.3917662287774708 0.38746926954634064 0.3956273540796493
% exponential2:age:  0.7798201255389867 0.755404691819801 0.800713066123444
% exponential2:massformed:  7.991977877526011 7.965030650848698 8.01724122467827
% exponential2:metallicity:  0.7950959975147073 0.7625852511324137 0.8216850961450373
% exponential2:tau:  0.18720279306834686 0.18164479536176947 0.19150189919719407
% lognormal:fwhm:  2.9180226797288973 2.8074751974002363 2.973671600851894
% lognormal:massformed:  9.494357283541525 9.490189012780103 9.498529225209797
% lognormal:metallicity:  0.539496784813187 0.5322367259651692 0.5471412283318265
% lognormal:tmax:  5.655128142220473 5.55493652207212 5.768589713612181
% nebular:logU:  -2.1275395742143166 -2.1999086015424067 -2.05412149157361
% redshift:  0.06899591065475572 0.0689895023123985 0.06900265218471469

\begin{figure}[h]
  \includegraphics[width=\textwidth]{../pipes/plots/r0_priors/phil_model_2_sfh}
  \includegraphics[width=\textwidth]{../pipes/plots/r4_exponential_burst/phil_model_2_sfh}
  \caption{
  Blind fit model No. 2, priori SFH (top) and median posterior SFH (bottom)
  with $\chi^2=1.6250106349682143$. Standard deviations of the posterior SFH
  are given in grey.
  }
  \label{fig:phil_model_2_sfhs}
\end{figure}

\begin{figure}[h]
  \includegraphics[width=\textwidth]{phil_model_2_posterior}
  \caption{Posterior distributions for fitted parameters on model No. 2.}
  \label{fig:phil_model_2_posterior}
\end{figure}

Model No. 2 is an example case of a negative burst detection. This is an
important result to verify; in order for the fit results to be reliable when
fitting real data, it is essential that the chosen priors do not give false
positives. Fig. \ref{fig:phil_model_2_sfhs} shows the priori and posterior SFH
for this model galaxy. Both SFHs feature an old, tau model component, with
excellent agreement on the age and mass formed by the component. The priori
component has $T_0=\SI{5.41}{\giga\year}$, $M=10^{10.2}M_\odot$, and
$\tau=1.5$, while the posterior component has $T_0=\SI{4.51}{\giga\year}$,
$M=10^{10.24}M_\odot$, and $\tau=1.71$. The small discrepancies in age and
decay time of this component can be explained by their flatter posterior
distributions (see below).

Perhaps more importantly, though, is that the priori SFH does not feature a
burst, and consequently, the posterior burst mass is very low, a mere
$10^{1.27}M_\odot$. What is more, the burst component is distributed over a
long period of time, with a relatively long decay time of $\tau =
\SI{1.71}{\giga\year}$, such that its contribution the SFR is negligible,
regardless of its age. This is the condition for a negative burst detection.

The posterior distribution of these key parameters (given in Fig.
\ref{fig:phil_model_2_posterior}) are not as tightly peaked, but this is not
completely unexpected. As discussed in \S\ref{sec:stellar_population_fitting},
the temporal resolution on stellar ages from old components is poor, so the
older component can vary somewhat in time and still produce a very similar
stellar population in the present day. Similarly, the burst is of sufficiently
low mass that it can vary freely across the parameter space with negligible
impact on the stellar population.

% parameter     median     16th percentile     84th percentile
% dust:Av:  0.8792126644473133 0.8703367567342202 0.8876472731201768
% exponential1:age:  7.908389955568364 7.603070847623988 8.457634683635277
% exponential1:massformed:  10.237047704928862 10.226823593039494 10.249145008568764
% exponential1:metallicity:  0.00919945903316217 0.008598217781738802 0.00975688301601339
% exponential1:tau:  1.709323179375498 1.63457413276435 1.852904109971942
% exponential2:age:  1.7778673467922534 0.6212267395127304 2.922064662802084
% exponential2:massformed:  3.4562035783022034 1.2727179039301946 5.805111536696616
% exponential2:metallicity:  1.3205964706480928 0.5124411732750468 2.086675458569264
% exponential2:tau:  1.0464626635822452 0.44128531892414374 1.6892799958378664
% nebular:logU:  -2.3065723113618564 -2.564054488793322 -2.093119902676116
% redshift:  0.08000488688448999 0.0799522285838497 0.08005315348948328

The fit on model No. 5 is an example case of how these priors still miss low
mass priori bursts. The priori model in Fig. \ref{fig:phil_model_5_sfhs}
features a double power law old component and young burst component of mass
$M=10^{7.7}M_\odot$ at $T_0=\SI{12.26}{\giga\year}$, but the corresponding
posterior component forms just $10^{3.32}M\odot$---not enough to affect the SFR
when distributed over a long period of time. The posterior distribution of the
mass formed by the burst, given in Fig. \ref{fig:phil_model_5_posterior}, is
wide and flat. In the wings of this distribution, there are some values which
approach the priori value, but at very low frequency. The fit, therefore,
doesn't find these values to be very likely. Additionally, the chi-squared
statistic for the posterior spectrum is quite good with respect to many of the
fitted results, with $\chi^2 \approx 0.306$, so despite the burst being several
orders of magnitude smaller in the posterior, this is still an very plausible
way to reconstruct the stellar population.

In summary, R4 set of priors reliably detects bursts of moderate to high mass,
as long as they occur in the most recent \SI{3.5}{\giga\year} of the model
galaxy's history. Crucially, the priors are not so tightly constrained that
they force false positives in the posterior distribution. The trade off for
this requirement is that sometimes the misses bursts of moderate age
(\SIrange{3.5}{7.5}{\giga\year} or low mass. The stellar population necessary
to account for stars produced in these bursts tend to be smeared into the
decaying old component. Even in fits which do positively detect bursts, the
mass of the burst tends to be underestimated, so this is clearly a limitation
of fitting with these priors. Again, this might be fixed by more tightly
constraining the age and decay times of the old component, but at the risk of
over fitting.

The chi-squared values for the posterior spectra are generally bad, with a few
exceptions. This is commonly the case even when the posterior SFH accurately
recreates the shape of the priori one. I hypothesize that this is due to a
systematic underestimation of the \textit{total} stellar mass formed, i.e. the
shape of the posterior SFH and spectrum is good, but total mass and flux values
are systematically lower than the priori galaxy. Further investigation is
needed to determine whether this is the case, but since the blue end of the
spectrum is more dramatically affected by this when you consider the age-
weighted mass formed, it might also explain why small mass bursts are missed.

\begin{figure}[h]
  \includegraphics[width=\textwidth]{phil_model_5_posterior}
  \caption{Posterior distributions for fitted parameters on model No. 5.}
  \label{fig:phil_model_5_posterior}
\end{figure}

\begin{figure}[h]
  \includegraphics[width=\textwidth]{../pipes/plots/r0_priors/phil_model_5_sfh}
  \includegraphics[width=\textwidth]{../pipes/plots/r4_delayed_burst/phil_model_5_sfh}
  \caption{
  Blind fit model No. 5, priori SFH (top) and median posterior SFH (bottom)
  with $\chi^2=0.30643246331094176$. Standard deviations of the posterior SFH
  are given in grey.
  }
  \label{fig:phil_model_5_sfhs}
\end{figure}

% parameter     median     16th percentile     84th percentile
% delayed:age:  8.83501662832162 8.519820505033731 9.168968916353453
% delayed:massformed:  10.397956599938258 10.391339420235703 10.405133311180329
% delayed:metallicity:  0.3132807043896933 0.30771134504388914 0.31873504377393214
% delayed:tau:  0.9949668644739154 0.9563350860520694 1.0356189144811663
% dust:Av:  0.392099804098769 0.3839782551743006 0.39866729226807895
% exponential2:age:  1.822742887145689 0.6174900074042676 2.9370305006919284
% exponential2:massformed:  3.3223708535326084 1.2262215632311155 5.754874771543047
% exponential2:metallicity:  1.249151161876338 0.4039612072741856 2.0513635481916204
% exponential2:tau:  0.9937853304518411 0.4295827267748906 1.6693788024708547
% nebular:logU:  -2.192720747989025 -2.2997584449121176 -2.092070491669038
% redshift:  0.030002834938625843 0.029983698303874765 0.030019661000926533

% \newpage
% \begin{figure}[t]
%   \includegraphics[width=\textwidth]{../pipes/plots/r0_priors/phil_model_10_sfh}
%   \includegraphics[width=\textwidth]{../pipes/plots/r4_lognormal_burst/phil_model_10_sfh}
%   \caption{Blind fit model No. 10, priori SFH (top) and median posterior SFH (bottom) with $\chi^2=15.943261900215944$. Standard deviations of the posterior SFH are given in grey.}
%   \label{fig:phil_model_10_sfhs}
% \end{figure}
%
% \begin{figure}[t]
%   \includegraphics[width=\textwidth]{phil_model_10_posterior}
%   \caption{}
%   \label{fig:phil_model_10_posterior}
% \end{figure}

% \newpage
% \begin{figure}[t]
%   \includegraphics[width=\textwidth]{../pipes/plots/r0_priors/phil_model_1_sfh}
%   \includegraphics[width=\textwidth]{../pipes/plots/r4_lognormal_burst/phil_model_1_sfh}
%   \caption{Blind fit model No. 1, priori SFH (top) and median posterior SFH (bottom) with $\chi^2=0.16108464297132902$.}
%   \label{}
% \end{figure}
% \lipsum[1]
% \begin{figure}[t]
%   \includegraphics{phil_model_1_posterior}
%   \caption{}
%   \label{}
% \end{figure}
%
% \newpage
% \begin{figure}[t]
%   \includegraphics[width=\textwidth]{../pipes/plots/r0_priors/phil_model_3_sfh}
%   \includegraphics[width=\textwidth]{../pipes/plots/r4_lognormal_burst/phil_model_3_sfh}
%   \caption{Blind fit model No. 3, priori SFH (top) and median posterior SFH (bottom) with $\chi^2=17250.67010982654$.}
%   \label{}
% \end{figure}
% \lipsum[1]
% \begin{figure}[t]
%   \includegraphics[width=\textwidth]{phil_model_3_posterior}
%   \caption{}
%   \label{}
% \end{figure}
%
% \newpage
% \begin{figure}[t]
%   \includegraphics[width=\textwidth]{../pipes/plots/r0_priors/phil_model_6_sfh}
%   \includegraphics[width=\textwidth]{../pipes/plots/r4_lognormal_burst/phil_model_6_sfh}
%   \caption{Blind fit model No. 6, priori SFH (top) and median posterior SFH (bottom) with $\chi^2=1502.0759895971512$.}
%   \label{}
% \end{figure}
% \lipsum[1]
% \begin{figure}[t]
%   \includegraphics[width=\textwidth]{phil_model_6_posterior}
%   \caption{}
%   \label{}
% \end{figure}
%
% \newpage
% \begin{figure}[t]
%   \includegraphics[width=\textwidth]{../pipes/plots/r0_priors/phil_model_7_sfh}
%   \includegraphics[width=\textwidth]{../pipes/plots/r4_lognormal_burst/phil_model_7_sfh}
%   \caption{Blind fit model No. 7, priori SFH (top) and median posterior SFH (bottom) with $\chi^2=1.0572371706120804$.}
%   \label{}
% \end{figure}
% \lipsum[1]
% \begin{figure}[t]
%   \includegraphics[width=\textwidth]{phil_model_7_posterior}
%   \caption{}
%   \label{}
% \end{figure}
%
% \newpage
% \begin{figure}[t]
%   \includegraphics[width=\textwidth]{../pipes/plots/r0_priors/phil_model_8_sfh}
%   \includegraphics[width=\textwidth]{../pipes/plots/r4_dblplaw_burst/phil_model_8_sfh}
%   \caption{Blind fit model No. 8, priori SFH (top) and median posterior SFH (bottom) with $\chi^2=4.205928604767932$.}
%   \label{}
% \end{figure}
% \begin{figure}[t]
%   \lipsum[1]
%   \includegraphics[width=\textwidth]{phil_model_8_posterior}
%   \caption{}
%   \label{}
% \end{figure}
%
% \newpage
% \begin{figure}[t]
%   \includegraphics[width=\textwidth]{../pipes/plots/r0_priors/phil_model_9_sfh}
%   \includegraphics[width=\textwidth]{../pipes/plots/r4_lognormal_burst/phil_model_9_sfh}
%   \caption{Blind fit model No. 9, priori SFH (top) and median posterior SFH (bottom) with $\chi^2=98.48856771268424$.}
%   \label{}
% \end{figure}
% \begin{figure}[t]
%   \lipsum[1]
%   \includegraphics[width=\textwidth]{phil_model_9_posterior}
%   \caption{}
%   \label{}
% \end{figure}

\newpage
\section{TDE Fitting Results}\label{sec:tde_fitting}

Having established some confidence in the reliability of fitting with the R4
priors, I then turned \textsc{Bagpipes} on the eight TDE-host spectra obtained
from the XSHOOTER instrument. Here follows a discussion of the R4 priors
applied to these targets, their inferred star formation histories for the chi-
squared best fit spectra, and the posterior distributions of their parameters,
with comparison to established literature on the targets.

Since there are no a priori values for the posterior SFHs and spectra, we can
not know for certain whether each fit represents the ``true'' values. As such,
the essential results of these fits are somewhat qualitative---does the fit
detect a burst or not? Analysis will be limited to whether the fitted SFH is
physically reasonable and/or consistent with observation.

\newpage
\subsection{ASASSN-14li, $z=0.0206$}\label{sec:ASASSN-14li}
\begin{figure}[h]
  \centering
  \includegraphics[width=\textwidth]{ASASSN14li_spec}
  \includegraphics[width=\textwidth]{ASASSN14li_sfh}
  \caption{
  Observed (blue) and fitted (orange) spectrum for the host of ASASSN-14li, and
  posterior SFH, with fitted spectrum $\chi^2=74130.89457269038$.
  }
  \label{fig:ASASSN14li_sfh}
\end{figure}

ASASSN-14li is perhaps one of the best known and studied tidal disruption
events. The best-fit posterior SFH for this TDE host features a lognormal old
component of mass $10^{9.51}M_\odot$ and $t_{max}=\SI{5.98}{\giga\year}$. The
burst component in this fit has mass $10^{8.91}M_\odot$ and
$T_0=\SI{12.43}{\giga\year}$. This gives a peak burst SFR much higher than any
of the blind fit results. These rates of star formation are actually more
consistent with observed, actively starburst galaxies, where star formation
rates commonly reach $10^2M_\odot\mathrm{yr}^{-1}$.\cite{Veilleux_2008} In
combination with the much longer, low level formation from the old component,
this is a strongly positive result. This is consistent with the findings of
Holoien et al. 2016a that PGC 043234 is  post-starburst
galaxy.\cite{Holoien_2016a}. There is also evidence that PGC 043234 is remnant
of a recent merger of two galaxies, which at some time prior to the tidal
disruption event hosted a type-II Seyfert.\cite{Prieto_2016} So while the host
is squarely post-starburst, it is worth noting that PGC 043234 is not
necessarily morphologically representative of all TDE hosts.

% r4_lognormal_burst \\
% r4_exponential_burst :  74287.13491573761 \\
% r4_dblplaw_burst :  74351.92695430398 \\
% r4_delayed_burst :  75626.48765873472 \\
% r4_lognormal_burst :  74130.89457269038 \\

% parameter     median     16th percentile     84th percentile
% dust:Av:  0.05098263083545196 0.04813568471489504 0.054433754544529554
% exponential2:age:  0.7656130076255705 0.7638063004738741 0.768246242489601
% exponential2:massformed:  8.905383130429595 8.89423045801556 8.917198002263056
% exponential2:metallicity:  1.4578520911723836 1.4447910938033477 1.4687321246075862
% exponential2:tau:  0.11722684527694574 0.11674721121378671 0.11778854880903881
% lognormal:fwhm:  2.9730405593202853 2.9415978145157866 2.991284482286685
% lognormal:massformed:  9.50522537973984 9.496666457963427 9.516922814281564
% lognormal:metallicity:  1.0003549300272776 0.9993051456778591 1.0017089522979168
% lognormal:tmax:  5.980146193951674 5.952834565182948 5.9940612348657005
% nebular:logU:  -2.0027322007706 -2.006068949885689 -2.000740615457861
% redshift:  0.02054054646568449 0.02053818187474229 0.020542639808073374

\newpage
\subsection{ASASSN-15oi, $z=0.0484$}\label{sec:ASASSN-15oi}
\begin{figure}[h]
  \centering
  \includegraphics[width=\textwidth]{ASASSN15oi_spec}
  \includegraphics[width=\textwidth]{ASASSN15oi_sfh}
  \caption{
  Observed (blue) and fitted (orange) spectrum for the host of ASASSN-15oi, and
  posterior SFH, with fitted spectrum $\chi^2=44017.90295163308$.
  }
  \label{fig:ASASSN15oi_sfh}
\end{figure}

\begin{figure}[h!]
  \centering
  \includegraphics[width=\textwidth]{ASASSN15oi_posterior}
  \caption{
  Posterior distributions for fitted parameters on the host of ASASSN-15oi.
  }
  \label{fig:ASASSN15oi_posterior}
\end{figure}

The posterior SFH for the host of ASASSN-15oi features a lognormal old
component of mass $10^{6.38}M_\odot$ and $t_{max}=\SI{3.31}{\giga\year}$. The
burst component in this fit has mass $10^{9.76}M_\odot$ and
$T_0=\SI{10.86}{\giga\year}$. In this fit, it would seems that the old
component actually has the negligible stellar mass, such that it is hardly
visible in Fig. \ref{fig:ASASSN15oi_sfh}. Although the old stellar mass is
within the lower range of galactic masses at $z<0.1$ catalogued by Baldry et
al. 2012 and Wright et al. 2017, it would imply that \textgreater 95\% of the
galaxy's mass was formed in the last \SI{3}{\giga\year}. This might be
plausible given the observed rates of star formation in active starbursts, and
that the total stellar mass of the galaxy is relatively small, but still seems
less physically reasonable.

What seems more plausible is that the exaggerated errors (light blue in Fig.
\ref{fig:ASASSN14li_sfh}--\ref{fig:iPTF16fnl_sfh}) on Telluric lines in the
spectrum, which occur primarily in the red and near-IR, are actually forcing
\textsc{Bagpipes} to discount the need for an older stellar population in the
fit. Indeed, we see from the posterior distributions of the component masses
and ages in Fig. \ref{fig:ASASSN15oi_posterior} that although the burst
parameters are tightly peaked, the mass formed in the old component could
plausibly be off by up to two orders of magnitude.

% r4_lognormal_burst \\
% r4_exponential_burst :  44187.55786809588 \\
% r4_dblplaw_burst :  44017.974181013255 \\
% r4_delayed_burst :  44019.20014908343 \\
% r4_lognormal_burst :  44017.90295163308 \\

% parameter     median     16th percentile     84th percentile
% dust:Av:  0.07640810516199385 0.0711248303132836 0.08223583948312649
% exponential2:age:  1.9305243305555009 1.9129861339724492 1.9461095547783132
% exponential2:massformed:  9.757288894215353 9.727731902627404 9.793858136851625
% exponential2:metallicity:  1.8276790321305791 1.8098831699330342 1.8468612686322659
% exponential2:tau:  0.28529021116961406 0.2823482650365451 0.28833227015962626
% lognormal:fwhm:  1.5321069676177346 0.6619235371269788 2.368997828914443
% lognormal:massformed:  6.383038494077901 5.478305437041224 7.418689818759305
% lognormal:metallicity:  1.4283668132122467 0.6625154248720727 2.105054197765897
% lognormal:tmax:  3.311013973502932 1.92846099345681 4.948748802286782
% nebular:logU:  -2.0038949924434957 -2.009254540499334 -2.0012725298682756
% redshift:  0.048413855141233045 0.04840937704606611 0.048418249388757696

% Although the redshift was originally reported as $z=0.02$, later spectra showed that the redshift of the host galaxy is $z=0.0484$, which was used in this fit.\cite{Arcavi_2015, Holoien_2016b}

\newpage
\subsection{AT2018fyk, $z=0.06$}\label{sec:AT2018fyk}
\begin{figure}[h!]
  \centering
  \includegraphics[width=\textwidth]{AT2018fyk_spec}
  \includegraphics[width=\textwidth]{AT2018fyk_sfh}
  \caption{
  Observed (blue) and fitted (orange) spectrum for the host of AT2018fyk, and
  posterior SFH, with fitted spectrum $\chi^2=49140.629786828795$.
  }
  \label{fig:AT2018fyk_sfh}
\end{figure}

The posterior SFH for the host of AT2018fyk features a lognormal old component
of mass $10^{10.25}M_\odot$ and $t_{max}=\SI{2.63}{\giga\year}$. The burst
component has mass $10^{8.66}M_\odot$ and $T_0=\SI{11.61}{\giga\year}$. This is
a slightly weak burst detection. The burst parameters do little to raise the
SFR above the quiescent rate, so this result is not exactly a starburst, but
still indicative of some recent, low-level star forming in the host. It is
worth noting that other R4 fits on this host all found bursts with SFRs as high
as $10M_\odot\mathrm{yr}^{-1}$, and only marginally worse chi-squared values of
49299, 49431, and 49443. With closer inspection, this could still be a post-
starburst candidate.

% r4_lognormal_burst \\
% r4_exponential_burst :  49299.44176377121 \\
% r4_dblplaw_burst :  49431.22634286046 \\
% r4_delayed_burst :  49443.68985368903 \\
% r4_lognormal_burst :  49140.629786828795 \\

% parameter     median     16th percentile     84th percentile
% dust:Av:  0.05753374521223291 0.05332317005875875 0.06144948174691915
% exponential2:age:  1.054821856385797 1.0153159033125843 1.0948487619419625
% exponential2:massformed:  8.61543829238758 8.581153794313689 8.646416920233435
% exponential2:metallicity:  1.7185623146070892 1.6779751817490571 1.7618610170697333
% exponential2:tau:  0.21902110439298791 0.2119953786303253 0.22596620482374816
% lognormal:fwhm:  0.3950463960209462 0.1870792062649676 0.6143836926002431
% lognormal:massformed:  10.246996416055287 10.244427485669771 10.249581873125544
% lognormal:metallicity:  0.9185127466881238 0.9130106591648134 0.923834975274212
% lognormal:tmax:  2.6311248957266256 2.1829082591358238 3.039662863486083
% nebular:logU:  -2.5953817431009893 -2.620943303129826 -2.56954606943412
% redshift:  0.059796780748178326 0.0597929546859466 0.059801428833008166

\newpage
\subsection{AT2019ahk, $z=0.026211$}\label{sec:AT2019ahk}
\begin{figure}[h]
  \centering
  \includegraphics[width=\textwidth]{AT2019ahk_spec}
  \includegraphics[width=\textwidth]{AT2019ahk_sfh}
  \caption{
  Observed (blue) and fitted (orange) spectrum for the host of AT2019ahk, and
  posterior SFH, with fitted spectrum $\chi^2=528057.3656910276$.
  }
  \label{fig:AT2019ahk_sfh}
\end{figure}

The posterior SFH for the host of AT2019ahk is another positive starburst
detection, featuring a delayed model old component of mass of
$10^{9.94}M_\odot$ and $T_0=\SI{5.54}{\giga\year}$. The burst component in this
fit has mass $10^{9.62}M_\odot$ and $T_0=\SI{12.42}{\giga\year}$, with peak
star formation of almost $25M_\odot\mathrm{yr}^{-1}$, as show in Fig.
\ref{fig:AT2019ahk_sfh}. This is a very similar result to that of ASASSN-14li,
with low-level star formation throughout the galaxy's life, and recent burst
which far exceeds the quiescent rate.

% r4_delayed_burst \\
% r4_exponential_burst :  528418.8168629913 \\
% r4_dblplaw_burst :  528786.5108022485 \\
% r4_delayed_burst :  528057.3656910276 \\
% r4_lognormal_burst :  528718.3921879242 \\

% parameter     median     16th percentile     84th percentile
% delayed:age:  7.5201250540978695 7.50836026038715 7.541352486105134
% delayed:massformed:  9.94474314208303 9.942468624134674 9.946724383919934
% delayed:metallicity:  0.7861243947678844 0.779596629211021 0.7922840791311964
% delayed:tau:  1.5444576470759133 1.5345051045915186 1.5548671136349284
% dust:Av:  0.8003814075932016 0.7988843668262824 0.8022380119127801
% exponential2:age:  0.6296672082410848 0.6287181578091492 0.6306634349078205
% exponential2:massformed:  9.618503174344696 9.61719115383671 9.619933982077773
% exponential2:metallicity:  2.2154203790790055 2.205791379473096 2.222772985192332
% exponential2:tau:  0.1849201168438975 0.18436480208511624 0.18556498257450266
% nebular:logU:  -2.7486451799452487 -2.7504461419737396 -2.7470480931311716
% redshift:  0.0263449902681042 0.02634440532841207 0.02634546636075307

\newpage
\subsection{AT2019azh, $z=0.022$}\label{sec:AT2019azh}
\begin{figure}[h]
  \centering
  \includegraphics[width=\textwidth]{AT2019azh_spec}
  \includegraphics[width=\textwidth]{AT2019azh_sfh}
  \caption{
  Observed (blue) and fitted (orange) spectrum for the host of AT2019azh, and
  posterior SFH, with fitted spectrum $\chi^2=100932.18767646234$.
  }
  \label{fig:AT2019azh_sfh}
\end{figure}

The posterior SFH for the host of AT2019azh features a delayed model old
component of mass $10^{9.06}M_\odot$ and $T_0=\SI{4.83}{\giga\year}$. The burst
component in this fit has mass $10^{9.55}M_\odot$ and
$T_0=\SI{12.57}{\giga\year}$. This result is superficially similar to that for
ASASSN-15oi, but the old component mass is much more substantial with respect
to the burst mass, in this fit. The very small decay time
($\tau=\SI{0.1}{\giga\year}$) on the burst component compresses the formation
of its stellar mass into a very short period with a high rate of star
formation, while the long decay time ($\tau=\SI{1.95}{\giga\year}$) on the old
component spreads it over over a long period of time. As such, its contribution
to the SFR is not visible in Fig. \ref{fig:AT2019azh_sfh}, even though its
contribution to the total mass formed is nearly equivalent.

% r4_delayed_burst \\
% r4_exponential_burst :  101469.19197835973 \\
% r4_dblplaw_burst :  101465.11902798309 \\
% r4_delayed_burst :  100932.18767646234 \\
% r4_lognormal_burst :  101454.60689188624 \\

% parameter     median     16th percentile     84th percentile
% delayed:age:  8.365288720823267 8.146628631509092 8.504078840994238
% delayed:massformed:  9.062312571171674 9.04899142203608 9.07402608902414
% delayed:metallicity:  0.9059826139899827 0.8964075740642858 0.9153502909178379
% delayed:tau:  1.9461848760247946 1.8906559665304545 1.9792001939045862
% dust:Av:  0.0002085667113198859 8.141090037844153e-05 0.00042491806204481545
% exponential2:age:  0.6256410303486765 0.6246372639317016 0.6264560330756612
% exponential2:massformed:  9.552269664694704 9.54738161216225 9.556321936697792
% exponential2:metallicity:  1.831238276352788 1.8248298461608183 1.837958151275079
% exponential2:tau:  0.10012124111600171 0.10004557455174809 0.1002435560350593
% nebular:logU:  -2.0010545825330666 -2.0019392006418717 -2.000405568756773
% redshift:  0.02214114933980705 0.022139743885776265 0.02214243001024535

\newpage
\subsection{AT2019dsg, $z=0.0512$}\label{sec:AT2019dsg}
\begin{figure}[h]
  \centering
  \includegraphics[width=\textwidth]{AT2019dsg_spec}
  \includegraphics[width=\textwidth]{AT2019dsg_sfh}
  \caption{
  Observed (blue) and fitted (orange) spectrum for the host of AT2019dsg, and
  posterior SFH, with fitted spectrum $\chi^2=80006.09829536812$.
  }
  \label{fig:AT2019dsg_sfh}
\end{figure}

The posterior SFH for the host of AT2019dsg features a delayed model old
component of mass $10^{9.88}M_\odot$ and $T_0=\SI{5.26}{\giga\year}$. The burst
component in this fit has mass $10^{10.58}M_\odot$ and
$T_0=\SI{9.40}{\giga\year}$. This is yet another positive starburst detection,
which again distinguishes itself from a component of intransient, low-level
star formation, and which is rapidly quenched before the present day.

% r4_delayed_burst \\
% r4_exponential_burst :  82563.8151623096 \\
% r4_dblplaw_burst :  81848.69517784887 \\
% r4_delayed_burst :  80006.09829536812 \\
% r4_lognormal_burst :  82054.99541292594 \\

% parameter     median     16th percentile     84th percentile
% delayed:age:  7.523514666251668 7.5113513446542175 7.534214285466046
% delayed:massformed:  9.88164797207937 9.880029925315881 9.88298310254143
% delayed:metallicity:  0.019165091355210325 0.018987494981708777 0.019331738087476853
% delayed:tau:  1.941765119373449 1.9407908045348816 1.9429934376101752
% dust:Av:  0.954656899419104 0.9534572275913585 0.9562581655632909
% exponential2:age:  3.38949788139676 3.3890417331325615 3.3898940343969777
% exponential2:massformed:  10.580578332707972 10.580451738854553 10.580725954133566
% exponential2:metallicity:  2.3739030389391 2.3705711264768565 2.377041711794388
% exponential2:tau:  1.0735214808813265 1.072732769776482 1.0742792612701677
% nebular:logU:  -3.036106338905969 -3.0377177095105017 -3.0340971042670386
% redshift:  0.05027470126334388 0.050271774603406934 0.05027685854725509

\newpage
\subsection{AT2019qiz, $z=0.01513$}\label{sec:AT2019qiz}
\begin{figure}[h]
  \centering
  \includegraphics[width=\textwidth]{AT2019qiz_spec}
  \includegraphics[width=\textwidth]{AT2019qiz_sfh}
  \caption{
  Observed (blue) and fitted (orange) spectrum for the host of AT2019qiz, and
  posterior SFH, with fitted spectrum $\chi^2=87438.51785338041$.
  }
  \label{fig:AT2019qiz_sfh}
\end{figure}

The posterior SFH for the host of AT2019qiz features a delayed model old
component of mass $10^{9.07}M_\odot$ and $T_0=\SI{5.66}{\giga\year}$. The burst
component in this fit has mass $10^{9.98}M_\odot$ and
$T_0=\SI{12.24}{\giga\year}$. This is a very similar result to that of
AT2019azh, in which two components of nearly equal stellar mass form that mass
on very different timescales. In this case, burst component has decay time
$\tau=\SI{0.21}{\giga\year}$ while the old component has
$\tau=\SI{1.99}{\giga\year}$. Again, we find that this galaxy had rates of star
formation comparable to observed starburst galaxies as recently as
\SI{1}{\giga\year} ago, and that this period of formation was rapidly quenched,
making it categorically post-starburst.

% r4_delayed_burst \\
% r4_exponential_burst :  88338.7683232356 \\
% r4_dblplaw_burst :  88346.50883758818 \\
% r4_delayed_burst :  87438.51785338041 \\
% r4_lognormal_burst :  90793.43597018946 \\

% parameter     median     16th percentile     84th percentile
% delayed:age:  7.527944307544574 7.517450323388695 7.536757664276217
% delayed:massformed:  9.067008182869145 9.059883655433673 9.074125746411527
% delayed:metallicity:  0.20540796042726123 0.20250312107485774 0.21008123964348271
% delayed:tau:  1.987200311366002 1.9816271669088061 1.992818805542456
% dust:Av:  1.0516156351780843 1.0480391424492903 1.0556850610038455
% exponential2:age:  0.9477035624254169 0.9466150964478025 0.948795028192377
% exponential2:massformed:  9.984776656965815 9.984275923477085 9.985499168966193
% exponential2:metallicity:  2.498898855293873 2.498122318184857 2.499344707090905
% exponential2:tau:  0.20636035995101426 0.20484693213247868 0.20749545015779258
% nebular:logU:  -2.7181110224531704 -2.7297916442935475 -2.7051244640475733
% redshift:  0.015299738365046962 0.015296524619805317 0.015302877316981591

\newpage
\subsection{iPTF16fnl, $z=0.018$}\label{sec:iPTF16fnl}
\begin{figure}[h]
  \centering
  \includegraphics[width=\textwidth]{iPTF16fnl_spec}
  \includegraphics[width=\textwidth]{iPTF16fnl_sfh}
  \caption{
  Observed (blue) and fitted (orange) spectrum for the host of iPTF16fnl, and
  posterior SFH, with fitted spectrum $\chi^2=122797.56521477805$.
  }
  \label{fig:iPTF16fnl_sfh}
\end{figure}

The posterior SFH for the host of iPTF16fnl features a tau model old component
of mass $10^{9.71}M_\odot$ and $T_0=\SI{3.42}{\giga\year}$. The burst component
in this fit has mass $10^{9.22}M_\odot$ and $T_0=\SI{12.51}{\giga\year}$. This
is another fit which is categorically post-starburst, with an old component of
low-level star formation and recent, rapidly quenched burst
component.\footnote{Important to note that Blagorodnova et al. 2017 estimates
this redshift to be $z=0.016328$, although this fit was run with
$z=0.018$.\cite{Blagorodnova_2017} This should not adversely affect the fit,
except perhaps to shift the model components slightly further back in time.}

% r4_exponential_burst \\
% r4_exponential_burst :  122797.56521477805 \\
% r4_dblplaw_burst :  122891.9232224498 \\
% r4_delayed_burst :  122831.05847432956 \\
% r4_lognormal_burst :  122800.99023633232 \\

% parameter     median     16th percentile     84th percentile
% dust:Av:  0.034241149477772584 0.029573525977327077 0.038386279250728496
% exponential1:age:  9.770783291093373 9.529239664622446 10.228223801519178
% exponential1:massformed:  9.710565031627318 9.700276348271476 9.720388214344052
% exponential1:metallicity:  1.0018237209344536 1.0001394870964546 1.005150328379511
% exponential1:tau:  0.6944782374064762 0.595701355181577 0.8160523434150315
% exponential2:age:  0.6716612949999092 0.6692064240461999 0.6750821430183963
% exponential2:massformed:  9.22085393371363 9.21752857794731 9.224259391914437
% exponential2:metallicity:  1.4813260403932302 1.4646165241912175 1.498049325537425
% exponential2:tau:  0.10004451800343839 0.10001137263891743 0.10010273358171788
% nebular:logU:  -2.0662154496360143 -2.1025206939240184 -2.035347433513074
% redshift:  0.01700003500026237 0.01700001110653572 0.017000088645588392

\newpage
Considering the concerns raised in \S\ref{sec:blind_fitting_results} regarding
a bias for diminished burst masses and the smearing of burst populations into
the old component, it is striking that all the TDE host galaxies exhibit such
strong signatures of recent, rapidly quenched, star forming activity. We might
hypothesize that if these biases are at play, then the ``true'' SFR peaks of
the TDE hosts were on the order of $10^2$ to $10^3M_\odot\mathrm{yr}^{-1}$, as
are observed in present day starburst galaxies. What is more, their posterior
SFHs are even more consistent than the blind fits with the paradigm that galaxy
formation occurs around cosmic noon, and that galaxies with nuclear activity
have variable SFRs after that epoch. This seems to corroborate the hypothesis
the TDEs, like AGN, preferentially occur in galaxies which are post-starburst,
and the photometric studies to that effect.

This result comes with some caveats, however. The posterior distributions of
key parameters (component age and mass) are tightly peaked for all eight fits,
with the exception of ASASSN-15oi, but the chi-squared values for the fitted
spectra are clearly abysmal. This does not necessarily invalidate the fits, as
the blind fitted spectra featured very poor chi-squared statistics even when
their posterior parameters accurately reproduced the priori SFH. Although the
posterior distributions for the nebular ionisation parameter all strongly
favour the high rates of ionisation associated with recent star formation,
\textsc{Bagpipes} does not seem to try very hard to fit the nebular emission
lines discussed in \S\ref{sec:starburst_agn_connection}, as is evident from
Figs. \ref{fig:ASASSN14li_sfh}--\ref{fig:iPTF16fnl_sfh}. \textsc{Bagpipes} does
not fit the observational spectrum with a high order polynomial as might be
assumed, but with the summation of emission and absorption functions given by
Eq. \ref{eq:luminosity_per_rframe_wavelength}, so it is had to understand why
this should be the case. Possibly the prior distribution needs to permit even
higher rates of ionisation, allowing $\log U \to 0$, in order to fit these
lines. The posterior distributions on this parameter are also tightly peaked,
though, suggesting this is not the case unless another maximum exists outside
the prior. The alternative interpretation would seem to be that galactic
spectra, being the composite of so many sources, are simply so noisy that we
must tolerate bad fits in order to understand the mechanisms underlying their
statistical behaviour.

\section{Conclusion}\label{sec:conclusion}

In summary, I have utilised the \textsc{Bagpipes} Python module to investigate
the prior assumptions necessary to detect recent ($<\SI{3.5}{\giga\year}$)
starburst activity in otherwise quiescent galaxies. I utilise literature values
on the statistical distributions of various galactic properties, including
stellar mass, population age, velocity dispersion, nebular ionisation, and dust
extinction, and the evolution of these properties over varying redshifts to
inform these assumptions. By simulating galactic spectra with known,
parametrised SFHs and blind fitting them, I determine that modest assumptions
about the form of the SFH must be made in order to detect these starbursts
reliably. Most important are the assumptions that the bulk of the stellar mass
in the galaxy is formed sometime in the first \SI{6.5}{\giga\year} since the
Big Bang, and that bursts older than \SI{3.5}{\giga\year}---the typical
lifetime of an F-type star---are too old to detect. By insisting on an old
component of star formation but allowing the burst mass to go to zero, I
differentiate between two possible fit results: the case of a single stellar
population (the null result), or two (a positive burst detection). Blind
fitting also illuminates the limitations of these priors: a tendency to
underestimate the burst mass and the total stellar mass formed, and a tendency
to smear stellar mass from the burst component into the quenching epoch of the
old component. Despite these drawbacks, the priors reliably reconstruct the
priori SFH of mock galaxies, including the presence of moderate to high mass
starbursts.

I then run fits with the same priors on eight TDE host spectra to investigate
whether they share a statistical property of AGN: that they occur
preferentially in galaxies which are post-starburst. The posterior SFHs for the
eight targets all feature strong starburst components, with most in the last
\SI{2}{\giga\year}, although not quite with the SFRs which are observed in some
active starburst galaxies. This result is consistent with previous photometric
analyses of the targets, which also find very young and blue stellar
populations. The fits also strongly affirm the paradigm that old star formation
peaks very early in cosmic history, $1.5<z<3$, despite priors which allow the
old component to occur much later. The fitted spectra for these targets do not
fit the observed spectra well, but there are good reasons to believe that the
reconstructed SFH is still valid. Finer tuning of the prior distributions on of
the nebular ionisation parameter in a later undertaking may improve this fit.

These results demonstrate that the methods applied by Carnall et al. 2019a to
quenched quiescent galaxies can also be used to infer star formation histories
containing recent variability, given reasonable assumptions based on the
statistical properties of similar galaxies at low redshifts. In the future, the
blind fitting method could be used to further investigate correlated parameters
in inferred star formation histories, providing additional insight to
complicated problem of nuclear/starburst driven feedback and quenching. A more
finely tuned set of priors turned on a larger sample of AGN and TDE hosts,
cross referenced by AGN-type, Hubble morphology, and other properties, could
well illuminate statistical trends in the SFHs of these taxonomies. Certainly
the SDSS archives of galactic spectra and photometry are a prime candidate for
this type of study.

\section{Acknowledgements}\label{sec:acknowledgements}

This work made use of the NASA/IPAC Extragalactic Database (NED), which is
operated by the Jet Propulsion Laboratory, California Institute of Technology,
under contract with the National Aeronautics and Space Administration. It also
made use of the The SAO/NASA Astrophysics Data System (ADS), which is operated
by the Smithsonian Astrophysical Observatory under NASA Cooperative Agreement.
Finally, this project made extensive use of the \textsc{AstroPy}
module.\cite{Astropy_2013, Astropy_2018} Special thanks to Andy Lawrence,
professor at the Royal Observatory Edinburgh, for his support and advice
throughout this project. Thanks also to Philip Short, a PhD candidate at the
University of Edinburgh, for his preparation of the XSHOOTER data and for
participating in my blind fitting trials. Thanks to Patrick Smith, for endowing
me with a love of empiricism and mentoring me through my formative school
years. And finally, thanks to my father, who even before I could stand, took me
for walks under the stars.

\printbibliography

\end{document}
